\scanpage{2}

Nous prendrons comme prétexte la théorie des fonctions $L$ pour développer plus avant le formalisme de la cohomologie étale. Trois parties

a) Continuer l'étude \inserted{de la coh. étale} avec des coefficients de torsion, en développant en particulier, dans ce cadre, \deleted{\unsure{plus}} \deleted{\unsure{uniquement}} avant le formalisme de dualité, \deleted{\unsure{donner}} donnant des formules de Künneth et de Lefschetz.

b) Étendre le formalisme à des coefficients \uline{$l$-adiques}.

c) Appliquer les résultats au formalisme des fonctions $L$.

1°) \uline{Site étale} de $X$ ;
\[
X_{\mathrm{ét}}:\quad\begin{array}{l}
\text{catégorie}\\
\text{familles couvrantes}\\
\text{faisceaux}\\
\text{faisceaux abéliens}\\
\text{cohomologie}
\end{array}
\]

2°) \uline{Exemple} $X-\unsure{\operatorname{Spec}}k$,
\[
X_{\mathrm{ét}}\simeq\widetilde X_{\mathrm{ét}}\simeq\operatorname{Ens}(\pi)
\qquad \pi=\operatorname{Gal}(\bar k/k)
\]

3°) $f:X\longrightarrow Y$ définit
\[X_{\mathrm{ét}}\;\unsure{-}\;Y_{\mathrm{ét}}\]
\[
f^*:\left\{\begin{array}{l}
Y_{\mathrm{ét}}\longrightarrow X_{\mathrm{ét}}\\
\widetilde Y_{\mathrm{ét}}\longrightarrow\widetilde X_{\mathrm{ét}}
\end{array}\right.
\]
\margin{commute aux $\unsure{\varinjlim}$ \unsure{quelc} et aux $\unsure{\varprojlim}$ finies}
\[f_*:\widetilde X_{\mathrm{ét}}\longrightarrow\widetilde Y_{\mathrm{ét}}\]
\margin{commute aux $\unsure{\varprojlim}$ quelc.}
\mspage{1}
